\chapter{结局回归与倾向性评分}

结局回归与各种形式的倾向性评分分析，是因果推断中最常用的参数方法。你或许会理所当然地感到奇怪：我们为什么拖了这么久，才拿出一章来讨论这些方法。到目前为止，我们讲的是 IP 加权、标准化与 g 估计——也就是 g 方法。把最常用的方法放在最不常用的方法之后来讲，在我们这边似乎是个古怪的选择。为什么不从更简单、也更广为使用的结局回归与倾向性评分方法讲起？因为这些方法一般并不成立。

更准确地说，本章无法尽数的那无数篇文献所描述的、较为简单的结局回归与倾向性评分方法，在较简单的设定下确实好用；但这些方法并不是为处理时变治疗下因果推断所伴随的复杂性而设计的。第三部分我们还会再讨论 g 方法，但对常规的结局回归与倾向性评分方法只会着墨较少。本章专门讨论那些常用、但对复杂纵向数据适用性有限的因果方法。\mnote{复杂纵向数据。}

\section{结局回归}

过去三章里，我们用 IP 加权、标准化与 g 估计来估计戒烟（处理）$A$ 对体重增加（结局）$Y$ 的平均因果效应。我们也说明了如何估计总体各子集内的平均因果效应：或者把分析限制在关心的子集上，或者在边缘结构模型（第 12 章）与结构嵌套模型（第 14 章）中加入乘积项。\mnote{提醒：我们把平均因果效应定义为 $\E[Y^{a=1,c=0}]-\E[Y^{a=0,c=0}]$。我们假定受治疗者与未受治疗者的可交换性，是在给定 $L$ 各变量——性别、年龄、种族、教育、吸烟强度与年限、日常生活中的体力活动、休闲锻炼以及体重——的条件下实现的。}以结构嵌套模型为例：这类模型含有处理 $A$ 与变量 $L$ 之间乘积项的参数，却不含 $L$ 本身的参数。这是结构嵌套模型一个吸引人的性质，因为我们关心的是在 $L$ 各水平内 $A$ 对 $Y$ 的因果效应，而不是 $L$ 与 $Y$ 之间的（非因果）关系。一种对 $L$–$Y$ 关系之函数形式不作任何假定的方法——结构嵌套模型的 g 估计——不会因这一关系被错误设定而产生偏倚。

反过来，如果我们愿意指定 $A$ 各水平内 $L$–$Y$ 的关联，就会考虑结构模型
\[
\E[Y^{a,c=0}\given L]=\beta_0+\beta_1 a+\beta_2 aL+\beta_3 L
\]
其中 $\beta_2$ 与 $\beta_3$ 是向量参数。戒烟 $A$ 对体重增加 $Y$ 在 $L$ 每一层内的平均因果效应，是 $\beta_1$ 与 $\beta_2$ 的函数；$L$ 每一层内在未治疗下的平均反事实结局，是 $\beta_0$ 与 $\beta_3$ 的函数。\marginnote{第 12 章中我们把这一模型称为\emph{伪边缘结构模型}：它具有边缘结构模型的形式，但估计其参数并不需要 IP 加权。稳定化 IP 权重 $SW^A(L)$ 全等于 1，因为该模型以整个向量 $L$ 为条件，而不是以 $L$ 的某个子集 $V$ 为条件。}[0.5cm]参数 $\beta_3$ 通常被称为 $L$ 的主效应，但用“效应”一词会引起误解，因为 $\beta_3$ 未必能解释为 $L$ 的因果效应（$L$ 可能存在混杂）。正如从我们的结构模型所见，$\beta_3$ 只是刻画了反事实 $Y^{a=0,c=0}$ 的均值如何随 $L$ 变化。关于像 $\beta_0$ 与 $\beta_3$ 这样没有因果解释的参数，见精要 15.1。

\begin{finepoint}[冗余参数]
设我们的目标是估计因果参数 $\beta_1$ 与 $\beta_2$。若通过拟合结局回归模型 $\E[Y^{a,c=0}\given L]=\beta_0+\beta_1a+\beta_2aL+\beta_3L$ 来估计，那么一般只有当 $\beta_0+\beta_3L$ 正确刻画了均值 $\E[Y^{a=0,c=0}\given L]$ 对 $L$ 的依赖时，$\beta_1$ 与 $\beta_2$ 的估计才是一致的。我们把参数 $\beta_0$ 与 $\beta_3$ 称为\emph{冗余参数}，因为它们不是我们主要关心的参数。

另一方面，若通过结构嵌套模型 $\E[Y^{a,c=0}-Y^{a=0,c=0}\given L]=\beta_1a+\beta_2aL$ 的 g 估计来估计 $\beta_1$ 与 $\beta_2$，那么一般只有当给定 $L$ 时治疗的条件概率 $\Pr[A=1\given L]$ 正确时，$\beta_1$ 与 $\beta_2$ 的估计才一致。也就是说，治疗模型的参数，例如 $\operatorname{logit}\Pr[A=1\given L]=\alpha_0+\alpha_1L$ 中的参数，现在成了冗余参数。

举例来说，如果某个协变量 $L$ 本应以线性项与二次项 $\beta_3L+\beta_4L^2$ 建模，却只用了线性项 $\beta_3L$，结局回归模型就会产生偏倚。结构嵌套模型不受结局回归模型错误设定的影响，因为结构模型并未指定 $L$–$Y$ 关系。然而，如果在治疗模型中错误设定了 $L$–$A$ 关系，用结构嵌套模型的 g 估计就会产生偏倚。对称地，结局回归模型也不受治疗模型错误设定的影响。对不随时间变化的固定治疗而言，选用哪种方法，归结为判断哪一组冗余参数——结局模型中的还是治疗模型中的——能被更准确地估计。可能的话，更好的选择是使用双重稳健方法（见精要 13.2）。
\end{finepoint}

若在 $L$ 的第 $l$ 层中所有人都接受治疗且都未被删失，反事实平均结局 $\E[Y^{a=1,c=0}\given L=l]$ 在可交换性、正值性与干预定义良好之下，等于未删失的受治疗者中相应的平均结局 $\E[Y\given A=1,C=0,L=l]$；未治疗者亦然。\mnote{$\beta_0$ 与 $\beta_3$ 刻画了 $Y^{a=0,c=0}$ 对 $L$ 的依赖关系；当模型被用来估计 (i) 平均反事实结局，以及 (ii) $L$ 各水平内的条件效应（在乘法尺度而非加法尺度上）时，这种依赖关系是必需的。}因此，上述结构模型的参数可以通过拟合结局回归模型
\[
\E[Y\given A,C=0,L]=\alpha_0+\alpha_1A+\alpha_2AL+\alpha_3L
\]
用普通最小二乘来估计，如 13.2 节所述。结局回归与第 3 章的分层一样，通过在 $L$ 的每一层内估计治疗的因果效应来调整混杂。如果变量 $L$ 足以调整混杂（与选择偏倚），且结局模型被正确设定，就不需要再做别的调整。也就是说，回归模型的参数 $\alpha$ 等于结构模型的参数 $\beta$。

在 13.2 节，结局回归只是通向标准化结局均值估计的中间一步；这里的结局回归则是整个过程的终点。我们不再把条件均值的估计标准化以得到边缘均值，而是直接比较条件均值的估计。在 13.2 节，我们拟合的回归模型在 $\beta_2$ 中只含一个乘积项（$A$ 与吸烟强度之间），也就是说，我们在模型中事先把大多数乘积项设为零。使用与 13.2 节相同的模型，这里得到的参数估计为 $\hat\beta_1=2.6$、$\hat\beta_2=0.05$。例如，对每天吸 5 支烟者，效应估计 $\widehat{\E}[Y\given A=1,C=0,L]-\widehat{\E}[Y\given A=0,C=0,L]$ 为 2.8（95\% 置信区间：1.5, 4.1）；对每天吸 40 支烟者为 4.4（95\% 置信区间：2.8, 6.1）。

结局回归的一种常见做法，是假定 $L$ 中任何变量都不修饰效应。此时模型不含任何乘积项，$\hat\beta_1$ 既是对治疗的条件平均因果效应的估计，也是对治疗的边缘平均因果效应的估计。在我们的例子里，不含任何乘积项的模型给出的估计为 3.5（95\% 置信区间：2.6, 4.3）kg。\mnote{代码：程序 15.1}

本章不必解释如何拟合结局回归模型，因为第 13 章估计参数化 g 公式的各个成分时已经做过了。用结局回归处理离散结局同样直截了当。对二值结局 $Y$，可以拟合 $\Pr[Y=1\given A=a,C=0,L]$ 的 logistic 模型。

\section{倾向性评分}

使用 IP 加权（第 12 章）与 g 估计（第 14 章）时，我们估计了每个个体在给定协变量 $L$ 下接受治疗的概率 $\Pr[A=1\given L]$。我们把这个条件概率记作 $\pi(L)$。对接受治疗概率低的个体，$\pi(L)$ 接近于 0；对接受治疗概率高的个体，$\pi(L)$ 接近于 1。也就是说，$\pi(L)$ 度量了在协变量 $L$ 所提供的信息之下，个体接受治疗的倾向。难怪 $\pi(L)$ 被称为\emph{倾向性评分}。

在一个理想的随机化实验中，一半个体被分配接受治疗 $A=1$，此时所有个体的倾向性评分都是 $\pi(L)=0.5$。还要注意，无论 $L$ 如何取，都有 $\pi(L)=0.5$。相反，在观察性研究中，有些个体可能比另一些更可能接受治疗。由于治疗分配不在研究者的控制之下，真实的倾向性评分 $\pi(L)$ 未知，因而必须从数据中估计。

在我们的例子里，可以对戒烟 $A$ 在给定协变量 $L$ 下的概率拟合 logistic 模型，以估计倾向性评分 $\pi(L)$。这与我们用于 IP 加权和 g 估计的是同一个模型。\mnote{代码：程序 15.2\\这里只考虑二值处理的倾向性评分。除 IP 加权、g 估计及其他相关的双重稳健估计量之外，倾向性评分方法很难推广到非二值处理。}

在该模型下，个体 22941 的估计倾向性评分最低（0.053），个体 24949 的最高（0.793）。图 15.1 给出了戒烟者 $A=1$（下）与未戒烟者 $A=0$（上）中估计倾向性评分的分布。\marginnote{图 15.1}[2.4cm]正如所料，戒烟者平均而言有更高的估计戒烟概率（0.312），未戒烟者为 0.245。如果 $\pi(L)$ 的分布在受治疗者 $A=1$ 与未受治疗者 $A=0$ 中相同，那么就不存在由 $L$ 引起的混杂，也就是说，因果图上不存在从 $L$ 到 $A$ 的开通路。

\begin{figure}[htbp]
\centering
\begin{tikzpicture}[x=1cm,y=1cm]
  \def\bw{0.72}
  \def\off{0.12}
  \def\sc{0.0175}
  % 上侧：A=0（未戒烟者）
  \foreach \v [count=\i from 0] in {13,123,192,258,203,172,98,61,41,23,9,5,3,0,0}{%
    \pgfmathsetmacro{\x}{\off+\i*\bw}%
    \pgfmathsetmacro{\h}{\v*\sc}%
    \pgfmathsetmacro{\xm}{\x+(\bw-0.1)/2}%
    \fill[white] (\x,0) rectangle (\x+\bw-0.1,\h);
    \draw[black!60,line width=0.35pt] (\x,0) rectangle (\x+\bw-0.1,\h);
    \node[above=0.6pt,inner sep=0pt,font=\fontsize{5.6}{6.2}\selectfont] at (\xm,\h) {\v};
  }
  % 下侧：A=1（戒烟者）
  \foreach \v [count=\i from 0] in {2,8,37,66,80,65,55,44,21,19,13,12,4,1,1}{%
    \pgfmathsetmacro{\x}{\off+\i*\bw}%
    \pgfmathsetmacro{\h}{-\v*\sc}%
    \pgfmathsetmacro{\xm}{\x+(\bw-0.1)/2}%
    \fill[black!78] (\x,0) rectangle (\x+\bw-0.1,\h);
    \node[below=0.6pt,inner sep=0pt,font=\fontsize{5.6}{6.2}\selectfont] at (\xm,\h) {\v};
  }
  % 坐标轴
  \draw[black!70,line width=0.4pt] (-0.05,0) -- (15*\bw+\off+0.02,0);
  \draw[black!70,line width=0.4pt] (-0.05,-1.95) -- (15*\bw+\off+0.02,-1.95);
  \draw[black!70,line width=0.4pt] (0,-1.95) -- (0,4.8);
  \foreach \y/\lab in {3.50/200,1.75/100,0/0,-1.75/100}{%
    \draw (-0.05,\y) -- (0,\y);
    \node[left=1.5pt,inner sep=1pt,font=\fontsize{5.8}{6.4}\selectfont] at (-0.05,\y) {\lab};
  }
  \foreach \lab [count=\i from 0] in {0.05,0.1,0.15,0.2,0.25,0.3,0.35,0.4,0.45,0.5,0.55,0.6,0.65,0.7,0.75,0.8}{%
    \pgfmathsetmacro{\x}{\off+\i*\bw}%
    \draw (\x,-1.95) -- (\x,-2.03);
    \node[below=1.5pt,inner sep=0.3pt,font=\fontsize{5.4}{6}\selectfont] at (\x,-2.03) {\lab};
  }
  \node[rotate=90,font=\fontsize{6.4}{7}\selectfont] at (-0.95,1.2) {例数};
  \node[below=12pt,font=\fontsize{6.4}{7}\selectfont] at (6.0,-2.03) {估计的倾向性评分};
  \node[right,font=\fontsize{5.8}{6.4}\selectfont] at (15*\bw+\off+0.05,4.55) {$A=0$};
  \node[right,font=\fontsize{5.8}{6.4}\selectfont] at (15*\bw+\off+0.05,-1.62) {$A=1$};
\end{tikzpicture}
\caption{}
\label{fig:15.1}
\end{figure}

倾向性评分 $\pi(L)$ 相同的个体，其某些协变量 $L$ 的取值通常不同。例如，两个 $\pi(L)=0.2$ 的个体可能在吸烟强度与锻炼上不同，但给定 $L$ 中全部变量，他们戒烟的可能性相同。也就是说，两人最终进入治疗组 $A=1$ 的条件概率相同。若考虑超总体中所有取某一给定 $\pi(L)$ 值的个体，这一组人会包含 $L$ 取不同值（例如吸烟强度与锻炼取不同值）的个体，但 $L$ 的分布在受治疗者与未受治疗者中相同，即 $A\indep L\given\pi(L)$。我们说倾向性评分在受治疗者与未受治疗者之间\emph{平衡}了协变量。

当然，倾向性评分只能平衡\emph{已测量}的协变量 $L$，这并不能防止未测量因素造成的残余混杂。随机化同时平衡已测量与未测量的协变量，因而是消除混杂的首选方法。\mnote{在随机化实验中，估计的 $\pi(L)$ 同时校正协变量的系统性与随机性失衡，因而优于对忽略随机失衡的真实 $\pi(L)$ 的调整。}关于平衡评分的正式定义，见技术注记 15.1。

\begin{techpoint}[平衡评分与预后评分]
如正文所述，倾向性评分 $\pi(L)$ 在受治疗者与未受治疗者之间平衡协变量。事实上，倾向性评分 $\pi(L)$ 是平衡评分最简单的例子。更一般地，\emph{平衡评分} $b(L)$ 是协变量 $L$ 的任一函数，只要
\[
A\indep L\given b(L).
\]
也就是说，对平衡评分的每个取值，协变量 $L$ 的分布在受治疗者与未受治疗者中相同。Rosenbaum and Rubin (1983) 证明：基于变量 $L$ 的可交换性与正值性，蕴含基于平衡评分 $b(L)$ 的可交换性与正值性。如果调整 $L$ 已足够，那么调整平衡评分 $b(L)$（包括倾向性评分 $\pi(L)$）也就足够。图 15.2 的因果图描绘了图 7.1 所示设定下的倾向性评分：可以把 $\pi(L)$ 看作位于 $L$ 与 $A$ 之间的中间节点，从 $L$ 到 $\pi(L)$ 有一条确定性的箭头。注意到 $\pi(L)$ 阻断了从 $A$ 到 $Y$ 的所有后门通路，我们也就给出了调整 $\pi(L)$ 之充分性的证明。

平衡评分 $b(L)$ 的一种替代是\emph{预后评分} $s(L)$，即协变量 $L$ 的一个函数，满足
\[
Y^{a=0}\indep L\given s(L).
\]
调整方法可以针对平衡评分与预后评分分别发展，但预后评分的方法需要更强的假定，且不易推广到时变治疗。关于预后评分的讨论，见 Hansen (2008) 与 Abadie et al. (2013)。
\end{techpoint}

像我们讨论过的所有因果推断方法一样，使用倾向性评分方法也需要可交换性、正值性与一致性这些识别条件。倾向性评分方法的正当性来自下面这个关键结果：受治疗者与未受治疗者在协变量 $L$ 各水平内的可交换性，蕴含二者在倾向性评分 $\pi(L)$ 各水平内的可交换性。也就是说，条件可交换性 $Y^a\indep A\given L$ 蕴含 $Y^a\indep A\given\pi(L)$。\mnote{若 $L$ 足以调整混杂与选择偏倚，那么 $\pi(L)$ 也足够。这一结果由 Rosenbaum 与 Rubin 在 1983 年发表的一篇开创性论文中导出。}此外，倾向性评分 $\pi(L)$ 各水平内的正值性——即没有个体的倾向性评分等于 1 或 0——成立，当且仅当第 2 章所定义的协变量 $L$ 各水平内的正值性成立。

在倾向性评分 $\pi(L)$ 各水平内的可交换性与正值性之下，可以用倾向性评分通过分层（包括结局回归）、标准化与匹配来估计因果效应。下面两节说明如何实施这些方法。作为第一步，我们必须先从观察数据估计倾向性评分 $\pi(L)$，再以估计的倾向性评分代替协变量 $L$ 去做分层、标准化或匹配。

\section{倾向性评分分层与标准化}

在识别条件之下，倾向性评分 $\pi(L)$ 取某一特定值 $s$ 的个体的平均因果效应，即
\[
\E[Y^{a=1,c=0}\given\pi(L)=s]-\E[Y^{a=0,c=0}\given\pi(L)=s],
\]
等于
\[
\E[Y\given A=1,C=0,\pi(L)=s]-\E[Y\given A=0,C=0,\pi(L)=s].
\]
这一条件效应可以通过把分析限制在真实倾向性评分取值为 $s$ 的个体上来估计。然而，倾向性评分 $\pi(L)$ 一般是可以在 0 与 1 之间取任意值的连续变量，因此两个个体不大可能恰好有相同的取值 $s$。例如，只有个体 22005 的估计 $\pi(L)$ 为 0.6563，这意味着我们无法通过比较该特定取值下的受治疗者与未受治疗者，来估计 $\pi(L)=0.6563$ 者的因果效应。

处理连续倾向性评分的一种办法，是划分出 $\pi(L)$ 取值相近（但不相同）的个体所构成的层。估计 $\pi(L)$ 的十分位数是常用的选择：把总体中的个体分入大小大致相等的 10 层，再在每一层内估计因果效应。\mnote{代码：程序 15.3}在我们的例子里，每个十分位约含 162 人。戒烟对体重增加的效应在各十分位之间从 0.0 到 6.6 kg 不等，但这些点估计周围的 95\% 置信区间很宽。

我们也可以拟合 $\E[Y\given A,C=0,\pi(L)]$ 的结局回归模型来得到这些效应估计：模型以治疗 $A$、估计 $\pi(L)$ 十分位的 9 个指示变量（其中一个十分位为参照水平，已包含在模型的截距中），以及 $A$ 与这些指示变量之间的 9 个乘积项为协变量。用估计 $\pi(L)$ 的十分位做结局回归的多数应用并不包含乘积项，即假定不存在 $\pi(L)$ 的效应修饰。在我们的例子里，不含乘积项的模型给出的效应估计为 3.5 kg（95\% 置信区间：2.6, 4.4）。关于倾向性评分的效应修饰，见精要 15.2。

以十分位或其他倾向性评分函数分层，会带来一个潜在问题：一般而言，连续 $\pi(L)$ 的分布在某些层（例如十分位）内会在受治疗者与未受治疗者之间不同。例如，若在某些层内受治疗者的平均 $\pi(L)$ 大于未受治疗者，那么在这些层内受治疗者与未受治疗者可能不可交换。前面几章用倾向性评分的函数，通过 IP 加权与 g 估计来估计结构模型的参数时，并没有这个问题，因为那些方法使用的是估计概率的数值，而不是十分位这样的分类变换。同样，把估计的倾向性评分 $\pi(L)$ 作为连续协变量（而不是一组指示变量）放入 $\E[Y\given A,C=0,\pi(L)]$ 的结局回归时，也没有这个问题。在我们例子里使用后一种做法，效应估计为 3.6（95\% 置信区间：2.7, 4.5）kg。

我们推断的有效性，取决于 $\pi(L)$ 与结局均值 $Y$ 之间关系的正确设定（我们假定为线性）。然而，由于倾向性评分是多维 $L$ 的一维概括，要防止这一关系被错误设定并不难：对倾向性评分改用三次样条等灵活模型，而不是单一线性项。\mnote{注意：二值处理 $A$ 的 IP 加权，其分母不是倾向性评分 $\pi(L)$，而是 $\pi(L)$ 的函数。对受治疗者（$A=1$）分母为 $\pi(L)$，对未受治疗者（$A=0$）为 $1-\pi(L)$。}注意，IP 加权与 g 估计对倾向性评分和结局之间的关系不作任何假定。

当 $\E[Y\given A,C=0,\pi(L)]$ 的参数假定正确，且可交换性与正值性成立时，该模型估计的是倾向性评分所有水平 $s$ 内的平均因果效应 $\E[Y^{a=1,c=0}\given\pi(L)=s]-\E[Y^{a=0,c=0}\given\pi(L)=s]$。如果我们关心的是整个研究总体中的平均因果效应 $\E[Y^{a=1,c=0}]-\E[Y^{a=0,c=0}]$，就要用倾向性评分的分布对条件均值 $\E[Y\given A,C=0,\pi(L)]$ 作标准化。其做法与第 13 章所述连续变量的做法相同，只是把变量 $L$ 换成估计的 $\pi(L)$。\mnote{代码：程序 15.4}注意，这一做法可以自然地在结局模型中纳入治疗 $A$ 与估计 $\pi(L)$ 之间的乘积项。在我们的例子里，标准化后的效应估计为 3.6（95\% 置信区间：2.7, 4.6）kg。

\section{倾向性评分匹配}

以倾向性评分 $\pi(L)$ 匹配的过程，类似于以单个连续变量 $L$ 匹配（见第 4 章）。倾向性匹配有许多形式，它们都试图构造一个匹配总体，使受治疗者与未受治疗者因 $\pi(L)$ 的分布相同而可交换。例如，可以把未受治疗者匹配到受治疗者：每个受治疗个体与一个（或多个）倾向性评分取值相同的未受治疗个体配对。由这些受治疗—未受治疗对（或组）所构成的、原总体的子集，就是\emph{匹配总体}。在给定 $\pi(L)$ 的可交换性与正值性之下，匹配总体中的关联度量是效应度量的一致估计：匹配总体中的关联风险比一致地估计匹配总体中的因果风险比。\mnote{匹配之后，匹配总体具有受治疗者的 $\pi(L)$ 分布、未受治疗者的 $\pi(L)$ 分布，或任何其他任意分布。}

同样，两个个体不大可能恰好有相同的倾向性评分 $\pi(L)$ 取值。在我们的例子里，倾向性评分匹配的做法是：为每个受治疗个体找出一个（或多个）$\pi(L)$ 取值接近的未受治疗个体。常见的做法是，把估计 $\pi(L)$ 取值为 $s$ 的受治疗个体，与取值在 $s\pm 0.05$（或其他某个小差值）之内的未受治疗个体匹配。例如，受治疗个体 1089（估计 $\pi(L)$ 为 0.6563）可以与未受治疗个体 1088（估计 $\pi(L)$ 为 0.6579）匹配。\mnote{匹配过去的一个缺点是，谁也不知道如何计算效应估计的方差。多亏 Abadie and Imbens (2006) 的工作，这已不再是问题。}定义“接近”的方式有很多，详细描述这些定义超出了本书的范围。

在倾向性匹配中定义“接近”，涉及偏倚—方差的权衡。如果接近的标准太宽松，$\pi(L)$ 取值相差较大的个体就会彼此匹配，匹配总体中 $\pi(L)$ 的分布在受治疗者与未受治疗者之间就会不同，可交换性便不成立。反之，如果接近的标准太严格，许多个体被匹配程序排除，就会得到近似的可交换性，但效应估计的 95\% 置信区间可能更宽。

“接近”的定义还与正值性的定义有关。在我们的戒烟例子里，估计 $\pi(L)$ 的分布在受治疗者与未受治疗者中于大部分取值范围上重叠（见图 15.1）。只有 2 名受治疗个体（占研究总体的 0.01\%）的取值高于任何未受治疗个体。上一节对估计 $\pi(L)$ 做结局回归时，我们实际上假定：缺少 $\pi(L)$ 估计值高的未受治疗个体是偶然所致——即随机性非正值——因而把所有个体都纳入分析。相反，多数倾向性匹配分析会认为这 2 名受治疗个体与任何未受治疗个体都不够接近，从而把它们排除。匹配并不区分随机性非正值与结构性非正值。\mnote{记住：正值性现在定义在倾向性评分的各水平内，即对所有满足 $\Pr[\pi(L)=s]$ 非零的 $s$，都有 $\Pr[A=a\given\pi(L)=s]>0$。}

上述讨论说明，匹配总体可能与目标（超）总体相差很大。理论上，倾向性匹配可以用来估计一个刻画良好的目标总体中的因果效应。例如，把每个受治疗个体与一个或多个未受治疗个体匹配、并排除未匹配的未受治疗者时，估计的是受治疗者中的效应（见精要 15.2）。但在实践中，倾向性匹配可能给出的是研究总体中一个难以描述的子系统内的效应估计。例如，在某种“接近”定义下，有些受治疗个体无法与任何未受治疗个体匹配，因而被排除在分析之外。结果，效应估计所对应的总体子集，是由那些匹配成功的估计倾向性评分取值所定义的。

倾向性匹配迫使研究者把分析限制在估计倾向性评分分布有重叠的治疗组，这一点常被说成该方法的优点。毕竟，谁也不想因为违反正值性而得到有偏的估计，对吧？然而，姑且不论与随机变异有关的问题（见上），基于倾向性评分的限制是有代价的。设想我们看过图 15.1 后得出结论：只能估计倾向性评分小于 0.67 者的戒烟效应。这些人是谁？并不清楚，因为个体并不会把倾向性评分纹在额头上。由于匹配总体刻画不良，很难评估效应估计向其他总体的可移植性。\mnote{即便每个受试者都把倾向性评分纹在额头上，总体仍可能刻画不良，因为相同的倾向性评分取值在不同情境下可能含义不同。}

当出现正值性方面的顾虑时，基于真实世界变量（例如年龄、吸烟支数）的限制会给出对因果效应更自然的刻画。在我们的戒烟例子里，估计 $\pi(L)>0.67$ 的这 2 名受治疗者，是研究中仅有的两位年过 50 且吸烟不足 10 年的人。我们可以把他们排除，并说明效应估计只适用于 50 岁以下的吸烟者，以及 50 岁及以上、吸烟至少 10 年的吸烟者。这样定义目标总体，比把它定义为“估计 $\pi(L)<0.67$ 者”更自然。

用倾向性评分来识别受治疗者与未受治疗者的重叠范围可能有用，但仅把研究总体限制在该范围之内，是一种偷懒的正值性保证方式。倾向性匹配自动获得的正值性，需要与“仅依据估计倾向性评分取值作限制时难以评估可移植性”这一点相权衡。

\begin{finepoint}[效应修饰与倾向性评分]
匹配估计与未匹配估计可能不同，一个原因是效应修饰。举个例子，考虑未受治疗者人数远多于受治疗者这一常见情形。倾向性匹配往往使几乎所有受治疗者都得到匹配，而许多未受治疗者未获匹配、因而被排除在分析之外。此时，匹配总体中因果效应修饰变量的分布将与受治疗者中相似。因此，匹配总体中的效应会更接近受治疗者中的效应，而不是用全体总体数据的方法所估计的效应。关于通过 IP 加权与标准化估计受治疗者中效应的其他办法，见技术注记 4.1。

倾向性评分各层之间的效应修饰，可以被解释为决策者知道自己在做什么的证据，例如医生倾向于治疗更可能从治疗中获益的病人（Kurth et al 2006）。然而，$\pi(L)$ 的效应修饰可能使估计的解释复杂化。设想一种定性效应修饰的情形：“医生，根据我们的研究，这种药对就诊时倾向性评分在 0.11 与 0.93 之间的病人有益，但可能使倾向性评分低于 0.11 的病人丧命”；或者“部长女士，我们只把这项教育干预用于倾向性评分低于 0.57 的儿童吧。”上述说法几乎没有政策意义，因为正如正文所论，它们不是用已测量的变量 $L$ 来表达的。

最后，除效应修饰之外，匹配估计与总体效应估计不同还有其他原因：未匹配者中正值性被违反、某个未测量混杂因子在匹配总体中比在未匹配总体中更常见（或更少见）、或者测量得更好（或更差），等等。正如第 4 章针对单个变量 $L$ 所论，表面上的效应修饰也可能由各倾向性评分层之间残余混杂的差异来解释。
\end{finepoint}

\begin{figure}[htbp]
\centering
\begin{tikzpicture}[causal,font=\small]
  \node (L) at (0,0) {$L$};
  \node (P) at (2.1,0) {$\pi(L)$};
  \node (A) at (4.5,0) {$A$};
  \node (Y) at (6.2,0) {$Y$};
  \draw[ed] (L) -- (P);
  \draw[ed] (P) -- (A);
  \draw[ed] (A) -- (Y);
  \draw[ed] (L.north) to[out=62,in=118] (Y.north);
\end{tikzpicture}
\caption{}
\label{fig:15.2}
\end{figure}

\section{倾向性评分模型、结构模型、预测模型}

本书第二部分介绍了因果推断的两类不同模型：倾向性模型与结构模型。现在我们来比较它们。

倾向性模型是给定变量 $L$ 时治疗 $A$ 之概率的模型，用来实现条件可交换性。本章我们用倾向性模型做匹配与分层，第 12 章用它做 IP 加权，第 14 章用它做 g 估计。\mnote{虽然倾向性评分是一维的，我们仍须用一个把治疗对高维 $L$ 作回归的模型来估计它。IP 加权与 g 估计同样如此。}倾向性模型的参数是冗余参数（见精要 15.1），没有因果解释，因为变量 $L$ 与治疗 $A$ 可能因许多原因而相关——并不只是因为 $L$ 导致 $A$。例如，在图 7.1 下可以把 $L$ 与 $A$ 之间的关联解释为 $L$ 对 $A$ 的效应，但在图 7.2 与图 7.3 下则不能。然而倾向性模型对因果推断有用，常常作为估计结构模型参数的基础，正如本章与前几章所述。

结构模型描述治疗 $A$ 与反事实结局 $Y^a$ 分布的某个成分（例如均值）之间的关系，这种关系可以是边缘的，也可以是在变量 $L$ 各水平内的。对连续治疗，结构模型常被称为剂量—反应模型。结构模型中治疗的参数不是冗余参数：它们具有直接的因果解释，即不同治疗取值 $a$ 下结局之差。我们介绍了两类结构模型：边缘结构模型与结构嵌套模型。边缘结构模型包含治疗、可能修饰效应的变量 $V$、以及治疗与 $V$ 之间乘积项的参数。$V$ 的选择只反映研究者对效应修饰的实质兴趣（见 12.5 节）。若不含任何协变量 $V$，模型就是真正边缘的。若把全部变量 $L$ 都作为可能的效应修饰变量纳入，边缘结构模型就成为伪边缘结构模型。结构嵌套模型包含治疗、以及治疗 $A$ 与 $L$ 中所有效应修饰变量之间乘积项的参数。\mnote{关于结构嵌套模型与伪半参数边缘结构模型之间的关系及其他微妙之处，见精要 14.1。}

我们把结局回归作为估计伪边缘结构模型参数以进行因果推断的方法。然而，结局回归也被广泛用于纯粹预测（而非因果）的目的。例如，网络零售商使用复杂的结局回归模型来预测哪些顾客更可能购买其产品。目的不是判定你的年龄、性别、收入、籍贯与既往购买是否对你当前的购买有因果效应，而是找出更可能购买的顾客，以便对他们定向投放特定的营销方案。这一切关乎关联，而非因果。同样，医生使用基于结局回归的算法来识别有发生重病或死亡高风险的病人。这些预测模型的参数未必有因果解释，模型中所有协变量的地位相同，即不存在治疗变量 $A$ 与变量 $L$ 之分。\mnote{一项研究发现，Facebook 的“点赞”能预测性取向、政治观点与人格特质（Kosinski et al. 2013）。这纯粹是预测，未必是因果。}

结局回归兼用于因果推断方法与预测，这造成了许多误解。其中最重要的误解之一与变量选择程序有关。当兴趣完全在于结局预测时，研究者可能想选出任何作为协变量纳入模型后能提高预测能力的变量。许多著名的变量选择程序——例如前向选择、后向剔除、逐步选择——以及机器学习中较新的进展，都被用来增强预测。对关心预测的研究者来说，这些是有力的工具，尤其是在处理非常高维的数据时。

不幸的是，统计学课程与教科书并不总是把因果推断与预测清楚地区分开。结果，这些用于预测模型的变量选择程序常被套用到因果推断模型上。这种错配的一个可能后果，是把多余的——甚至有害的——协变量纳入倾向性模型与结构模型。具体而言，把预测算法用于因果推断模型，可能导致方差膨胀与偏倚增大。

问题之所以出现，是因为人们普遍（却错误地）相信：倾向性模型应当尽可能好地预测治疗 $A$。倾向性模型并不需要很好地预测治疗，它们只需要包含能保证可交换性的变量 $L$。与治疗强关联、但对保证可交换性并非必需的协变量，无助于减少偏倚。若把这些协变量纳入 $L$，调整可能得到方差极大的估计，甚至放大已有的偏倚。\mnote{倾向性分析报告 Mallows's $C_p$ 之类预测力度量并不罕见。这些度量对因果推断的意义值得怀疑。}

看下面这个例子。设某项研究中所有个体都在 Aceso 医院或 Panacea 医院就诊。Aceso 医院的医生给 99\% 的个体治疗 $A=1$，Panacea 医院的医生给 99\% 的个体 $A=0$。设变量“医院”对结局没有效应（除了通过其对治疗 $A$ 的效应），因而对实现条件可交换性并非必需。假设我们仍决定把“医院”作为协变量加入倾向性模型。目标总体中的倾向性评分 $\pi(L)$ 对 Aceso 医院的人约为 0.99，对 Panacea 医院的人约为 0.01；但出于偶然，我们得到的研究总体中，可能出现 Aceso 医院所有人在 $L$ 所定义的某些层内都有 $A=1$，或 Panacea 医院所有人都有 $A=0$ 的情形。也就是说，我们的效应估计可能有近乎无限的方差，而混杂一点也没有减少。治疗现在被预测得很好，这对因果推断而言无关紧要。\mnote{如果我们完美地预测了治疗，那么所有受治疗者都会有 $\pi(L)=1$，所有未受治疗者都会有 $\pi(L)=0$。二者将没有重叠，分析也就无法进行。}

除方差膨胀之外，以预测的态度为因果推断模型（倾向性模型与结局回归模型皆然）选择变量，还可能导致自找的偏倚。例如，把对撞节点作为协变量纳入，可能造成系统性偏倚，尽管对纯粹预测而言对撞节点可能是有效的协变量；把工具变量（见下一章）纳入，可能放大未测量变量所致的偏倚。我们将在第 18 章回到这些问题。

一般地，基于模型的因果推断方法——倾向性模型与结构模型——要求协变量的函数形式不被错误设定。为降低模型错误设定的可能，我们使用三次样条这样的灵活设定，而不是线性项。此外，这些因果推断方法要求可交换性、正值性与干预定义良好，才能得到无偏的因果推断。下一章我们介绍一类非常不同的因果推断方法，它不要求治疗的可交换性。
